% Appendix I: contamination diagnostics and the decontaminated condition (plan 0013 §4.3).
% Method from the implemented code; results from the ledger rows of 2026-10-03 to 2026-10-07 (all seven
% systems rebuilt at five cutoffs). Keys: results.tex, _build/fill_decon.py; DECON-* typed from the ledger.
\ifdefined\nwvtwo
\section{Contamination Diagnostics and the Treated Text}
\else
\section{Contamination Diagnostics and the Decontaminated Condition}
\fi
\label{app:contamination}

A generator may recognize published novels \citep{chang2023speak}. We compare aligned
original, names-only and renamed-and-paraphrased text, assessing recognition separately
from later-fact disclosure and cutoff compliance. Renaming follows prior practice
\citep{gupta2025whodunit, zhang2024infinitebench}; cross-condition differences motivate
contamination diagnostics \citep{dekoninck2024constat}. Because a treatment can damage
the narrative \citep{sun2025emperor}, we examine gold-fact fidelity and document errors.
The procedure aims to preserve facts, attribution, reveal order and cross-volume
coreference, but has no independent validation of complete fidelity. Reduced recognition
does not rule out memorization \citep{cheng2025contamination}.


\begin{figure}[t]
\centering
% F5 (one column, v2): gold recall of each system on the original and on the treated text,
% means over t=1..5. Dot positions are typed from the results.tex keys M1-anne-avg-<S> and
% M1-annedecon-avg-<S>; re-type them if those keys change. The printed changes are \res keys.
\begin{tikzpicture}[x=4.3cm, y=0.36cm, font=\figfont,
  orig/.style={circle, draw=nwgray, fill=white, line width=0.5pt, minimum size=5.2pt, inner sep=0pt},
  treat/.style={circle, fill=nwaccent, minimum size=3.2pt, inner sep=0pt},
  sysn/.style={anchor=west, font=\figfont\bfseries},
  dif/.style={anchor=east, font=\figfont, text=nwink}]
\foreach \x/\l in {0/0, 0.25/.25, 0.5/.50, 0.75/.75, 1/1}{
  \draw[nwgrid] (\x,0.55) -- (\x,-6.65);
  \node[text=nwgray] at (\x,-7.2) {\l};}
\node[anchor=west, font=\figfont\itshape, text=nwgray] at (-0.5,1.05) {system};
\node[font=\figfont\itshape, text=nwgray] at (0.5,1.05) {gold recall};
\node[anchor=east, font=\figfont\itshape, text=nwgray] at (1.24,1.05) {change};
% rows: name / y / original / treated / change
\newcommand{\dumb}[5]{%
  \node[sysn] at (-0.5,#2) {#1};
  \draw[nwgray, line width=0.7pt] (#3,#2) -- (#4,#2);
  \node[orig] at (#3,#2) {}; \node[treat] at (#4,#2) {};
  \node[dif] at (1.24,#2) {#5};}
\dumb{B1 long context}{0}{0.848}{0.845}{\res{M1-annedecon-diff-B1}}
\dumb{B2 prefix retrieval}{-1}{0.834}{0.784}{\res{M1-annedecon-diff-B2}}
\dumb{B3 LightRAG}{-2}{0.536}{0.470}{\res{M1-annedecon-diff-B3}}
\dumb{\sys}{-3}{0.802}{0.843}{\res{M1-annedecon-diff-B5}}
\dumb{B4 graph baseline}{-4}{0.761}{0.704}{\res{M1-annedecon-diff-B4}}
\draw[nwgrid, densely dotted] (-0.5,-4.55) -- (1.24,-4.55);
\dumb{X2 full text}{-5.1}{0.814}{0.744}{\res{M1-annedecon-diff-X2}}
\dumb{X1 closed book}{-6.1}{0.683}{0.003}{\best{\res{M1-annedecon-diff-X1}}}
% legend
\node[orig] at (-0.47,-8.25) {}; \node[anchor=west] at (-0.44,-8.25) {original text};
\node[treat] at (0.2,-8.25) {}; \node[anchor=west] at (0.23,-8.25) {treated (renamed, paraphrased)};
\end{tikzpicture}

\caption{Gold recall on original and treated text (means over $t=1..5$) and the change,
treated minus original. Without source text, recall collapses; with it, recall is largely
retained.}
\label{fig:deconrecall}
\end{figure}

\paragraph{Treatment.}
The treatment operates paragraph by paragraph, retaining identifiers, chapter boundaries
and scene breaks so citations remain aligned.

\emph{Step 1, names.} A reviewed map of \res{DECON-map} entries replaces character,
family and local place names with invented forms. Longer forms precede shorter ones
in whole-word replacement, with consistent full names, surnames and nicknames.
Mechanical checks require stand-ins to be absent from the source names and every part
of a mapped full name to be mapped. An initial map failed both checks after one volume:
a main character's stand-in matched two minor characters' first names, and a mapped
full name left its 148 bare-first-name occurrences unchanged. We rebuilt the map and
discarded that treated volume. Recurring minor characters and local places are renamed;
countries and distant cities remain. A surname also used as a compass direction is
replaced only in person-name contexts.

\emph{Step 2, paraphrase.} The generator rewrites each renamed paragraph under
instructions to preserve facts, events, numbers, places, ordering, names and dialogue
speakers, with no additions or summarization. Reapplying the map replaces any restored
source names. Nonempty rewrites between half and twice the source length are accepted;
otherwise the renamed paragraph is retained and counted as unparaphrased.
Residual source names are counted separately. In an initial treated volume,
approximately one paragraph in 25, containing one word in 12, retained at least four
fifths of its four-word sequences. We therefore request one retry when at least half
the source's four-word sequences remain. A second near-copy is retained but counted
as verbatim. Failed model calls stop the run, preventing silent fallback.

\emph{Step 3, gold.} The same map transforms the original gold facts and cast;
facts are neither added nor removed, and reveal volumes are unchanged.
All systems in Section~\ref{sec:setup} are rebuilt with unchanged code, configuration
and prompts. A names-only variant omits Step~2 and is used for recognition probes,
rather than complete system rebuilding.

\paragraph{What the condition can and cannot show.}
If systems mainly use supplied evidence, much of their original-text coverage may survive
the treatment. Reduced recognition may instead reduce closed-book coverage, where no
source text is supplied. These are diagnostic expectations, not an identification
strategy for each system's memory contribution: rewriting also changes wording,
entity resolution and sometimes meaning.

A prefix-bounded system can disclose later facts through extraction or inference errors,
generation or residual recognition. A treated-text disclosure by the full-text diagnostic
still demonstrates a failure to obey the cutoff while later evidence is available.
A model can recognize a plot without its original names, so the treatment cannot establish
the absence of memorization.

\paragraph{The treated text.}
The treated series contains \res{DECON-paras} paragraphs and
\res{DECON-words} words before and after treatment. Of these paragraphs,
\res{DECON-paraphrased} were paraphrased, with the remainder renamed only;
\res{DECON-verbatim} remained near-copies after retrying. Paraphrase lengthens the series
overall, although each accepted rewrite stays within the length bounds of Step~2.
Whole-name checks find no
mapped source name unchanged. A broader search finds
\res{DECON-residual-forms-n} inflected or derived source-name occurrences in
\res{DECON-residual-forms} paragraphs: chiefly family-name plurals (``the Gardners''),
plus a college-derived adjective and a sentence explaining a name's origin.
Recognition probes include these residual forms.

Some residuals reach generated pages. Across \res{DECON-pages-scanned} character pages
from seven systems, the name explained in the source sentence appears in
\res{DECON-jamesina-pages} pages of six systems. A family name retained in a plural
and a chapter title (``The Gardners' Call'') appears on
\res{DECON-gardner-pages} pages of four systems. Three other family names and the
college name, each from a residual form, appear on seven pages. A further
closed-book match is a common surname in an invented story.

Another source name has no literal residual: \res{DECON-pei-pages} pages from prefix
retrieval and the full-text diagnostic print ``Prince Edward Island'', absent from
the treated paragraphs and chapter titles. The renamed island's neighboring provinces
remain, permitting geographic inference. This output therefore does not distinguish
inference from residual recognition.

\paragraph{Fidelity.}
An unblinded LLM annotator checked all 111 gold facts against their treated evidence
paragraphs. It judged \res{DECON-fidelity} still expressed, one with altered wording:
a sleeve style is renamed in its evidence paragraph but given its source name in
the next paragraph. This asks whether the treated passage retains the fact.
The stricter paragraph-alone audit from Appendix~\ref{app:annotation}, yielding
\res{HUMR-gold} for the original text, was not repeated.

Errors also occur outside gold evidence. One paraphrase changes two children's ages
into an interval of minutes, after which prefix retrieval reports twins born six
minutes apart. The statement counts as supported against the supplied treated text.
Treatment errors found during analysis are documented here and left in place.

\paragraph{Recognition.}
On 200 aligned paragraphs spanning five volumes, the generator restores a masked
character name with accuracy \res{MEM-cloze-orig} on original text and
\res{MEM-cloze-decon} on treated text. Title identification is
\res{MEM-title-orig}, \res{MEM-title-remap} and \res{MEM-title-decon}
for original, names-only and combined treatment, respectively; author identification
is \res{MEM-author-orig}, \res{MEM-author-remap} and \res{MEM-author-decon}.
Renaming reduces title recognition more than author recognition. Adding paraphrase
reduces both further, to approximately one third of their names-only levels, but
roughly one treated paragraph in nine is still identified as belonging to the series.
Only recognition was probed for names-only text; complete systems were rebuilt on
original and combined-treatment text.

% The renamed-only control was probed for recognition; it is not an additional
% complete condition of the wiki comparison.
\begin{table}[t]
\centering\small
\setlength{\tabcolsep}{5pt}
\renewcommand{\arraystretch}{1.12}
\begin{tabular}{@{}lrr@{}}
\toprule
\nwhead{Text condition} & \nwhead{Title} & \nwhead{Author} \\
\midrule
Original & \res{MEM-title-orig} & \res{MEM-author-orig} \\
Names remapped & \res{MEM-title-remap} & \res{MEM-author-remap} \\
Remapped and paraphrased & \cellcolor{nwhighcyan}\res{MEM-title-decon} & \cellcolor{nwhighcyan}\res{MEM-author-decon} \\
\bottomrule
\end{tabular}
\caption{Recognition ablation on 200 aligned paragraphs. Values are the share for which
the frozen generator identifies an Anne title or its author. Renaming and paraphrasing
both reduce recognition, while residual recognition remains. Cyan highlights the combined
treatment. This probe does not establish
that a wiki fact was memorized.}
\label{tab:recognition}
\end{table}


\paragraph{Systems on the treated text.}
Table~\ref{tab:decon} compares original and treated systems;
Table~\ref{tab:decon-cells} reports every cutoff. Gold facts, reveal volumes and
labeling rules are transferred from the original; inverse name mapping allows
the same per-fact tests. Claude Opus~5.5 finalized the AI labels for all systems on both
texts, including precision and citation sufficiency.

Each condition has one build. Differences in entity resolution and paraphrase fidelity
preclude interpreting small cross-condition changes as isolated memory effects.
The closed-book diagnostic has a much larger recall reduction, from
\respt{M1-anne-avg-X1} to \respt{M1-annedecon-avg-X1}, with
\res{M1-annedecon-count-X1} matching fact--cutoff assessments remaining:
a death also present in its invented story. This is consistent with recognition
contributing substantially to original-text closed-book recall.

\paragraph{What the closed-book generator writes.}
Given stand-in names without text, the generator still writes full pages.
Of 200 sampled assertions, \res{M3-annedecon-unsup-X1} are unsupported.
The stand-in for a shy book-one farmer becomes a sea captain founding a cult;
a heroine's friend becomes a duelist at a magical academy; later cutoffs introduce
a street gang and a starship. Supported judgments concern generic statements
about gender, status, siblings or traits. No assertion carries a citation.

\paragraph{The disclosure by \sys.}
At $t=2$ and $t=3$, \sys lists both leads as spouses, a book-four fact
(\res{M4-annedecon-total-B5}, counting the marriage for both characters at both cutoffs).
These are two reciprocal subject facts from one prediction-to-marriage failure mechanism,
rather than four independent errors (Table~\ref{tab:disclosure-units}).
Another \res{M4-annedecon-uncertain-B5} verdicts are uncertain, involving the
same courtship line as in the original. The marriage citations point to a
book-two neighbor's prediction. Extraction stored a presumed marriage at
volume two, and rendering turned presumed spouses into spouses.
B4, using the same graph, prints ``anticipated future marriage'' at $t=2$,
labeled implied disclosure (\res{M4-annedecon-total-B4}).

The original build did not extract this relation from the corresponding passage.
The treated passage is cutoff-eligible, but the claim loses the distinction
between a character's prediction and an established event. Volume filtering
cannot repair that attribution and modality error.

\paragraph{Long-context prompting on the treated text.}
Mean recall is \respt{M1-annedecon-avg-B1} on treated text versus
\respt{M1-anne-avg-B1} on original text; precision is
\respt{M3-annedecon-avg-B1} versus \respt{M3-anne-avg-B1}.
Both treated estimates are lowest at $t=4$ (Table~\ref{tab:decon-cells}).
The close means do not establish equivalence from single builds.

Citation sufficiency is \respt{M8-annedecon-t1-B1} and
\respt{M8-annedecon-t2-B1} for one and two volumes, then falls to
\respt{M8-annedecon-t3-B1}, \respt{M8-annedecon-t4-B1} and
\respt{M8-annedecon-t5-B1}, versus \respt{M8-anne-t3-B1},
\respt{M8-anne-t4-B1} and \respt{M8-anne-t5-B1} on the original.
Most insufficient citations accompany supported assertions. Many retain the
paragraph number but lose the volume or chapter, pointing, for example, from
a volume-one fact to the corresponding identifier in volume three.
Such identifiers pass the harness's existence check.

The cause is untested. Familiarity with the original might help locate evidence,
or paraphrased passages might be harder to distinguish. Neither explanation
is isolated by this comparison.

\paragraph{The full-text diagnostic on the treated text.}
Despite the cutoff instruction, X2 discloses \res{M4-annedecon-total-X2}
assessed later facts on treated text, versus \res{M4-anne-total-X2} on the original.
Per-cutoff counts at $t=1..4$ are \res{M4-annedecon-t1-X2},
\res{M4-annedecon-t2-X2}, \res{M4-annedecon-t3-X2} and
\res{M4-annedecon-t4-X2}; original counts are \res{M4-anne-t1-X2},
\res{M4-anne-t2-X2}, \res{M4-anne-t3-X2} and \res{M4-anne-t4-X2}.
At $t=1$, the guardian's page describes the neighbor's widowhood and move,
the heroine's college attendance and her marriage ``years later''.
It cites volumes two to four, whose identifiers the harness accepts because
all volumes are supplied. Reduced recognition therefore does not make the
instruction reliable; observed disclosure is higher at the first cutoff.
Familiarity with the original's book boundaries is one untested explanation.

Treated recall is lower at every cutoff (\respt{M1-annedecon-avg-X2} versus
\respt{M1-anne-avg-X2}); mean precision is also lower
(\respt{M3-annedecon-avg-X2} versus \respt{M3-anne-avg-X2}).
Most unsupported statements are misattributed biographies, and citation errors
also omit volumes. The heroine's $t=1$ page returned malformed JSON three times,
closing a list with the wrong bracket. We corrected that bracket in the
recorded first response without another model call
(Appendix~\ref{app:reproducibility}).

\paragraph{Namesakes made by the name map.}
Whole-word renaming creates two collisions observed in the systems receiving the most
text, B1 and X2. First, one stand-in replaces a shared first name across characters,
making a book-one minor character a namesake of a schoolmate; passages about one
are then cited for the other. Second, a book-two twin's stand-in is a short form
of the stand-in for the heroine's romantic name, also borne by her friend's
daughter in book four. With book four in context, the twin's page describes
the friend (B1 at $t=4$; X2 at $t=2$, $4$ and $5$) or the heroine
(X2 at $t=3$), whereas original-text builds describe the twin.

This page contributes unsupported assertions and two stated X2 disclosures
at $t=3$. Neither mechanical map check detects these collisions. They are
treatment defects retained in the reported system scores.
